\section{Mercy in the heart, hope on the road}
Both readings begin from a gift and refuse to make its recipient
self-sufficient. Their decisive questions differ. The first asks
what the already-forgiven person owes a fellow servant. It follows
the movement from remission to responsibility: received mercy must
become inward mercy, and cruelty is a grave contradiction of the gift.
The second asks how the already-delivered people can still suffer.
It follows the movement from release to rest: divine help sustains
the journey without making every stage painless or every traveler
automatically faithful.

The difference is clearest at the Offertory. In the first reading
Job prevents an abusive conclusion: the Gospel's condemned servant
is not a model for diagnosing every wounded person. In the second
Job supplies the central difficulty and its principal evidence:
uprightness and wounds coexist, and Gregory distinguishes hostile
intention, permitted action, and divine guardianship. Conversely,
the Gospel is the first reading's center and the second reading's
necessary warning. Freedom received can be used against its giver.

The Communion prevents either reading from becoming sentimental.
The merciful person may still ask God to judge injustice. The
protected pilgrim may still ask when help will become manifest.
Augustine's whole-Christ interpretation and Bellarmine's longing
for final peace carry the future hope; Schuster's present judgment
in the loss of truth supplies a different emphasis. The difference
survives within their common refusal to treat wickedness as
inconsequential.

The prayers bring both movements to the same altar without making
them the same argument. Restoration is God's gift; immortal
nourishment must reach the interior life. The mission additions
extend that gift outward to all peoples, and Luke's admitted
commemoration gives cross-bearing service a place within the
household's intercession. Mercy must not close upon its recipient,
and hope must not stop at the first deliverance. The family asks
to live what it receives and to reach the life for which it is fed.

\clearpage
